% Preamble
\documentclass[11pt, aspectratio=169]{beamer}
	% Packages
	\usepackage[english]{babel}
	\usepackage{lipsum}
	\usepackage{mathptmx}
	\usepackage{xcolor}
	% Themes
	\usefonttheme[onlymath]{serif}
	\usetheme[nofirafonts]{focus}
	\renewcommand{\thempfootnote}{\arabic{mpfootnote}}
	% Cover
	\title{Microeconometrics Using Stata}
	\subtitle{Linear regression extensions: Exercises}
	\author{
		Jhon R. Ordoñez
		\footnote{\label{1} National University of San Cristóbal de Huamanga}
		\footnote{\label{2} Faculty of Economic, Administrative and Accounting Sciences}
		\footnote{\label{3} Professional School of Economics}
	}
	\date{\today}
% Body
\begin{document}
	% Cover
	\begin{frame}
		\maketitle
	\end{frame}
	% Outline
	\begin{frame}[t]
		\frametitle{Outline}
		\tableofcontents
	\end{frame}
	% Section 1
	\section{Exercise 1}
		% Slide 1.1
		\begin{frame}
			\frametitle{Exercise 1}
				\begin{block}{Exercise 1}
					Fit the model in \textcolor{red!50}{\textbf{section 4.2}} on levels, except use all observations
					rather than those with just positive expenditures, and report robust
					standard errors. Predict medical expenditures. Use \textcolor{blue}{\textbf{correlate}} to obtain
					the correlation coefficient between the actual and fitted value, and
					show that, upon squaring, this equals $R^{2}$. Show that for the linear
					model \textcolor{blue}{\textbf{margins}} with the \textcolor{blue}{\textbf{dydx(*)}} option reproduces the \textbf{OLS}
					coefficients. Now, use \textcolor{blue}{\textbf{margins}} with an appropriate option to obtain
					the average income elasticity of medical expenditures.
				\end{block}
		\end{frame}
	% Section 2
	\section{Exercise 2}
		% Slide 2.1
		\begin{frame}
			\frametitle{Exercise 2}
			\begin{block}{Exercise 2}
				Consider the example of \textcolor{red!50}{\textbf{section 4.2.4}} and the third method that
				computes predictions at \textcolor{blue}{\textbf{suppins==1}} for all observations and at
				\textcolor{blue}{\textbf{suppins==0}} and then compares the average predictions. The text did
				this following OLS regression in levels and obtained an estimated
				treatment effect of $\$725$. Repeat this example but with predictions
				from the log-linear model that correct for retransformation bias using
				Duan’s method. You should obtain an estimate of $\$2,048$.
			\end{block}
		\end{frame}
	% Section 3
	\section{Exercise 3}
		% Slide 3.1
		\begin{frame}
			\frametitle{Exercise 3}
			\begin{block}{Exercise 3}
				Fit the model in \textcolor{red!50}{\textbf{section 4.2}} on levels, using the first $2,000$
				observations. Use these estimates to predict medical expenditures for
				the remaining $1,064$ observations, and compare these with the actual
				values. Note that the model predicts very poorly in part because the
				data were ordered by \textcolor{blue}{\textbf{totexp}}.
			\end{block}
		\end{frame}
	% Section 4
	\section{Exercise 4}
		% Slide 4.1
		\begin{frame}
			\frametitle{Exercise 4}
			\begin{block}{Exercise 4}
				Reconsider the out-of-sample prediction analysis of \textcolor{red!50}{\textbf{section 4.3.}}
				Modify the \textbf{DGP} such that the regression intercept for the last $10$
				observations is $5$, whereas it is fixed at $1$ for the first $90$ observations.
				Generate $10$ out-of-sample predictions, and compare them with the
				case in which no such intercept shift occurs.
			\end{block}
		\end{frame}
	% Section 5
	\section{Exercise 5}
		% Slide 5.1
		\begin{frame}
			\frametitle{Exercise 5}
			\begin{block}{Exercise 5}
				Repeat the analysis of the previous question but under the assumption
				that the intercept remains unchanged in the prediction period but the
				slope parameter increases or decreases by $20\%$. Compare the results
				with the no-change scenario.
			\end{block}
		\end{frame}
	% Section 6
	\section{Exercise 6}
		% Slide 6.1
		\begin{frame}
			\frametitle{Exercise 6}
			\begin{block}{Exercise 6}
				Refer to the empirical example of \textcolor{blue}{\textbf{section 4.3.}} Define two group-specific
				indicator variables, \textcolor{blue}{\textbf{d1}} for the female group and \textcolor{blue}{\textbf{d2}} for the
				male group. Sort the sample observations by group; females are in group 1. 
				Using the full sample, run the earnings regression given in the section; include \textcolor{blue}{\textbf{d1}} and \textcolor{blue}{\textbf{d2}} group indicators, but exclude the intercept
				term. Test the hypothesis that the coefficients of \textcolor{blue}{\textbf{d1}} and \textcolor{blue}{\textbf{d2}} are equal.
			\end{block}
		\end{frame}
	% Section 7
	\section{Exercise 7}
		% Slide 7.1
		\begin{frame}
			\frametitle{Exercise 7}
			\begin{block}{Exercise 7}
				This is a continuation of the previous question. Replace regressors in
				the regression by new ones generated by multiplying the original
				regressors by \textcolor{blue}{\textbf{d1}} and \textcolor{blue}{\textbf{d2}}. 
				The new regressors are group specific.
				Regress the earnings variable on the new generated group-specific
				regressors. Use the \textcolor{blue}{\textbf{test}} command to test for the equality of the two
				sets of group-specific coefficients.
			\end{block}
		\end{frame}
	% References
	\begin{frame}[t,allowframebreaks]
		\frametitle{References}
		\bibliography{biblio.bib}
		% Style
		\bibliographystyle{apalike}
		% Authors
		\nocite{cameron2022-I}
		\nocite{cameron2022-II}
	\end{frame}
	% Cover
	\begin{frame}
		\maketitle
	\end{frame}
\end{document}